% Technical-supplement fragment for EviTrace.
% This file is intended to be included with \input{...}; it is not a
% standalone document.  Author-supplied metadata that is not recoverable
% from the frozen experiment artifacts is deliberately left as a visible
% placeholder.
\providecommand{\placeholder}[1]{\textbf{[TODO--AUTHOR: #1]}}

\section{Additional Experimental Details}
\label{sec:supp-experiments}

This supplement records the data transformations, evaluation rules,
controlled-analysis protocols, and artifact contracts used in the paper.
Conditioned on the frozen atomization and Evidence Map response caches,
downstream preprocessing is deterministic.  All reported verifier backbones
are adapted with LoRA.  The evidence scorer and verifier are trained
separately: the pairwise scorer objective uses map-induced preferences,
whereas the verifier uses label-token cross-entropy.

\subsection{Datasets, Versions, and Preprocessing}
\label{sec:supp-data}

\subsubsection{Dataset releases and splits}

Table~\ref{tab:supp-dataset-splits} gives the exact number of claims loaded
from each release.  LIAR-RAW and RAWFC associate each claim with raw,
claim-linked reports~\citep{Yang2022CofCED}.  SciFact uses the released
corpus of 5,183 scientific abstracts~\citep{Wadden2020SciFact}.

\begin{table}[t]
\centering
\scriptsize
\setlength{\tabcolsep}{2.5pt}
\begin{tabular}{lrrrr}
\toprule
\textbf{Dataset} & \textbf{Train} & \textbf{Val./Dev.}
& \textbf{Test} & \textbf{Total} \\
\midrule
LIAR-RAW & 10,065 & 1,274 & 1,251 & 12,590 \\
RAWFC    & 1,612  & 200   & 200   & 2,012 \\
SciFact  & 809    & 300   & 300   & 1,409 \\
\bottomrule
\end{tabular}
\caption{Claim counts in the files consumed by the pipeline.  The SciFact
test file contains claims but no public verdict or rationale labels.}
\label{tab:supp-dataset-splits}
\end{table}

\begin{table}[t]
\centering
\scriptsize
\setlength{\tabcolsep}{2pt}
\begin{tabular}{llrrr}
\toprule
\textbf{Dataset} & \textbf{Label} & \textbf{Train}
& \textbf{Val./Dev.} & \textbf{Test} \\
\midrule
LIAR-RAW & pants-fire  & 812  & 115 & 86 \\
         & false       & 1,958& 259 & 249 \\
         & barely-true & 1,611& 236 & 210 \\
         & half-true   & 2,087& 244 & 263 \\
         & mostly-true & 1,950& 251 & 238 \\
         & true        & 1,647& 169 & 205 \\
\midrule
RAWFC & false & 514 & 66 & 66 \\
      & half  & 537 & 67 & 67 \\
      & true  & 561 & 67 & 67 \\
\midrule
SciFact & support    & 332 & 124 & -- \\
        & contradict & 173 & 64  & -- \\
        & NEI        & 304 & 112 & -- \\
\bottomrule
\end{tabular}
\caption{Class counts after label normalization.  Dashes denote hidden
SciFact test labels.}
\label{tab:supp-class-counts}
\end{table}

The LIAR-RAW six-class order is
\(\{\textit{pants-fire},\textit{false},\textit{barely-true},
\textit{half-true},\textit{mostly-true},\textit{true}\}\).
RAWFC's source labels \(1/3/5\) are mapped to
\(\textit{false}/\textit{half}/\textit{true}\), respectively.
For SciFact, a claim with no annotated evidence in the labeled release is
assigned NEI; otherwise the released evidence stance is normalized to
\textit{support} or \textit{contradict}.  The prediction aliases supplied
to the verifier are A--F for LIAR-RAW and A--C for the three-class tasks.
The alias-to-label mapping is fixed before training and is inverted before
metric computation.

\begin{table*}[t]
\centering
\small
\setlength{\tabcolsep}{4pt}
\begin{tabularx}{\textwidth}{@{}lXXX@{}}
\toprule
\textbf{Dataset} & \textbf{Source and access}
& \textbf{Frozen release identifier} & \textbf{Integrity record} \\
\midrule
LIAR-RAW &
Official CofCED/LIAR-RAW release with claim-linked reports &
\placeholder{release date/version or archive identifier} &
\placeholder{SHA-256 for train/val/test and license/access note} \\
RAWFC &
Official CofCED/RAWFC release with Snopes-linked report text &
\placeholder{release date/version or archive identifier} &
\placeholder{SHA-256 for train/val/test and license/access note} \\
SciFact &
Official release containing \texttt{corpus.jsonl} and the three claim files &
\placeholder{release date and source archive URL} &
\placeholder{SHA-256 for the archive and four extracted files} \\
\bottomrule
\end{tabularx}
\caption{Dataset provenance fields to be frozen with the camera-ready
artifact.  Hashes refer to the unmodified downloaded files.}
\label{tab:supp-dataset-provenance}
\end{table*}

\subsubsection{Text normalization and report materialization}

All claims and report fields are converted to strings, non-breaking spaces
and line breaks are replaced by ordinary spaces, consecutive whitespace is
collapsed, and leading/trailing whitespace is removed.  Empty report
elements are discarded.  LIAR-RAW preserves each released
\texttt{event\_id}, report identifier, link, and domain.  For RAWFC, every
non-empty element of the released \texttt{evidence} list is materialized as
one report with a deterministic within-claim index.  SciFact abstracts are
indexed by the released document identifier and zero-based sentence
position; titles and sentence boundaries from the release are preserved.
Retrieval queries are formed from the normalized claim and its atom
renderings; verifier evidence is formed from the indexed report or abstract
fields described above.

For raw reports that are not already sentence-tokenized, sentence splitting
protects a fixed list of common abbreviations and splits at sentence-final
punctuation followed by a capital letter or digit.  Empty sentences and
sentences shorter than 10 characters are removed.  Each retained sentence
keeps its claim, report, and sentence identifiers, so every selected unit
can be traced back to its source.

\subsubsection{Claim-aware chunks and candidate pools}

The exact adjacent-boundary chunker, its token definition, and both resolved
parameterizations are specified in Section~\ref{supp:chunking}.  The released
manifest records the chunk text, sentence interval, anchor sentence, token
count, relevance, boundary scores, and resolved chunker configuration.

Multi-path retrieval issues one query for the complete claim and one query
for each atom.  Each atom route retains 20 candidates.  Candidate relevance
combines normalized dense, lexical-overlap, and BM25-like scores using
weights \(0.70/0.20/0.10\); MMR uses \(\lambda=0.70\).
The claim-route and atom-route results are deduplicated using the canonical
text key defined in Section~\ref{supp:retrieval} to form the Atom--Union
pool.  Ties are broken by the frozen candidate index.

\subsection{Evaluation Protocol}
\label{sec:supp-evaluation}

\subsubsection{Classification metrics}

Let \(C\) be the number of labels.  For class \(c\), precision, recall, and
F1 are
\begin{equation}
\begin{aligned}
P_c&=\frac{TP_c}{TP_c+FP_c},\\
R_c&=\frac{TP_c}{TP_c+FN_c},\\
F_{1,c}&=\frac{2P_cR_c}{P_c+R_c}.
\end{aligned}
\end{equation}
A ratio with a zero denominator is defined as zero.  Macro scores are the
unweighted arithmetic means over the fixed label inventory:
\begin{equation}
\begin{aligned}
P_{\mathrm{macro}}&=\frac{1}{C}\sum_c P_c,\\
R_{\mathrm{macro}}&=\frac{1}{C}\sum_c R_c,\\
F_{1,\mathrm{macro}}&=\frac{1}{C}\sum_c F_{1,c}.
\end{aligned}
\end{equation}
Thus Macro-F1 is not computed as the harmonic mean of
\(P_{\mathrm{macro}}\) and \(R_{\mathrm{macro}}\).
Accuracy is the number of correctly classified claims divided by the
number of evaluated claims.  LIAR-RAW and RAWFC results use the official
test examples and report percentages.

The verifier predicts through label-token scoring rather than unconstrained
free-form parsing.  At the answer position, the logit of each allowed
single-token alias is extracted, the highest-logit alias is selected, and
the frozen alias map converts it to the dataset label.  A missing or
non-finite alias logit is treated as an evaluation error and is recorded,
not silently removed.

For SciFact, the released prediction schema, evaluation command, and raw
evaluation output are preserved in the artifact bundle.  No SciFact headline
result is reproduced in this supplement.

\subsubsection{Training, selection, and repeated runs}

All reported verifier runs use LoRA.  The common and dataset-specific
optimization settings are given in
Table~\ref{tab:supp-lora-settings}; the released resolved configuration,
rather than a short experiment name, is the authoritative hyperparameter
record.  LIAR-RAW and RAWFC use their respective validation splits for
configuration and checkpoint choices, followed by one official-test
evaluation.

\placeholder{State the number of independent runs represented
by every main-table cell.  If a cell is a single canonical run, label it as
such; otherwise report the aggregation rule and dispersion.}

\subsubsection{Baseline provenance and comparison scope}

The primary LIAR-RAW/RAWFC comparison is restricted to systems that consume
the claim-linked reports at inference.  Systems using a corpus-wide
benchmark, multiple benchmarks, open-web search, or gold evidence are
retained only as broader-access context.  A published score is transcribed
without recomputation unless the corresponding row is explicitly marked as
a reproduction.  Table~\ref{tab:supp-baseline-provenance} is the
camera-ready provenance ledger; each entry must point to a concrete source
table and state whether the source reports one run, a mean, or a selected
best run.

\begin{table*}[t]
\centering
\scriptsize
\setlength{\tabcolsep}{3.2pt}
\begin{tabularx}{\textwidth}{@{}p{0.10\textwidth}
p{0.22\textwidth}YYY@{}}
\toprule
\textbf{Scope} & \textbf{Methods} & \textbf{Evidence/backbone}
& \textbf{Score provenance} & \textbf{Aggregation/comparability note} \\
\midrule
Paired reports &
dEFEND; SBERT-FC; GenFE; GenFE-MT &
\placeholder{input boundary and backbone/size per row} &
\placeholder{source paper, table, split, and metric per row} &
\placeholder{copied/reproduced; runs and best/mean/single} \\
Paired reports &
MUSER; CofCED &
\placeholder{input boundary and backbone/size per row} &
\placeholder{source paper or cited reproduction and table} &
\placeholder{runs and aggregation; note reproduction status} \\
Paired reports &
L-Defense; G-Defense; EVICheck; DelphiAgent &
\placeholder{LLMs, auxiliary components, and report access} &
\placeholder{source paper, exact variant, table, and split} &
\placeholder{runs and any additional supervision} \\
Broader access &
DeReC; FFRR; ReAct; Verify-and-Edit; RAV &
\placeholder{benchmark/web/gold evidence and backbone} &
\placeholder{source paper or cited reproduction and table} &
Context only; not ranked as a matched paired-report comparison \\
SciFact &
VERISCI; Paragraph-Joint; PrunE; QMUL-SDS; VerT5erini; ARSJoint; RerrFact &
\placeholder{corpus, retrieval universe, and backbone per row} &
\placeholder{source paper, table, development split, scorer} &
\placeholder{pipeline vs.\ evidence-provided setting and runs} \\
\bottomrule
\end{tabularx}
\caption{Baseline provenance template.  This table must be completed from
the cited publications before release; a method name alone is not a
reproducible score source.}
\label{tab:supp-baseline-provenance}
\end{table*}

\subsection{Controlled Analyses for RQ2}
\label{sec:supp-rq2}

All RQ2 analyses use LIAR-RAW and the Ministral-3-8B verifier unless a
cross-verifier comparison is explicitly named.  Within each ablation
family, the claim atoms, report chunks, evaluation examples, label mapping,
and LoRA training recipe are held fixed.  Variant-specific evidence traces
are materialized for train, validation, and test, and the verifier is
LoRA-tuned on the representation it receives at evaluation; the experiment
does not replace evidence only at test time.

\subsubsection{Evidence-organization variants}

For each claim \(i\), let \(K_i\) be the number of complete evidence units
visible under the full EviTrace condition.  Every evidence-bearing control
in Table~\ref{tab:supp-selector-mechanics} is capped at the same \(K_i\); if its
pool contains fewer than \(K_i\) units, it exposes the entire pool and logs
the underfill.  Count matching is performed after deduplication and before
prompt token guarding.

\begin{table*}[t]
\centering
\small
\setlength{\tabcolsep}{3.5pt}
\begin{tabularx}{\textwidth}{@{}l l l X@{}}
\toprule
\textbf{Variant} & \textbf{Candidate pool} & \textbf{Ranking/order}
& \textbf{Mechanical intervention} \\
\midrule
No evidence & none & none &
The verifier receives the claim and label instruction but no evidence
units. \\
Random & raw report pool & uniform random order &
Sample \(K_i\) units without replacement from the claim's raw evidence
pool and preserve the sampled order.
\placeholder{insert the frozen sampling seed.} \\
Claim route only & claim-route pool & retrieval score &
Remove every atom-route-only candidate; take the highest-scoring \(K_i\)
claim-query candidates. \\
Atom routes only & atom-route pool & retrieval score &
Remove the whole-claim route; deduplicate the union of atom-query routes
and take its highest-scoring \(K_i\) candidates. \\
Retrieval only & Atom--Union & retrieval score &
Use the full union pool but replace state-conditioned scoring with
descending frozen retrieval score. \\
Rule only & Atom--Union & map-priority trace &
At each prefix, choose the lexicographic
\(\Delta_{\mathrm{map}}\) winner using direct/non-direct state progress,
new-atom coverage, relation novelty, map quality, retrieval relevance,
and stable index; no learned \(S_\theta\) is used. \\
Shuffled trace & Atom--Union & random permutation &
Keep the evidence set and \(K_i\) produced by \(S_\theta\), then apply one
frozen random permutation to its order.
\placeholder{insert the shuffle seed and tie policy.} \\
EviTrace (Full) & Atom--Union & \(S_\theta\) trace &
Iteratively score remaining candidates, append the deterministic
argmax, update atom states, and stop under the full capacity policy. \\
\bottomrule
\end{tabularx}
\caption{Mechanical definitions for the evidence-organization ablation.
All evidence-bearing variants are matched to the full method's
claim-specific visible evidence count.}
\label{tab:supp-selector-mechanics}
\end{table*}

\subsubsection{Evidence Map field interventions}

The four map-field variants keep the same Atom--Union pool and
minmax\((5,10)\) capacity policy, alter only the named map input, regenerate
all dependent state features and preferences, and retrain the scorer on the
degraded representation.  For \emph{No relation},
\(r_{ij}\leftarrow\textit{background}\) for every pair.  For
\emph{No directness}, \(d_{ij}\leftarrow\textit{context}\).  For
\emph{Constant confidence}, \(c_{ij}\leftarrow1.0\).  For
\emph{Map disabled}, all map-derived features and state transitions are
disabled and retrieval-derived features remain available.  The full
condition retains the original LLM-generated
\((r_{ij},d_{ij},c_{ij})\).  The average visible \(K\) is reported because
these interventions may change whether the resolution-based stopping
condition fires.

\subsubsection{Balanced cross-verifier mechanism experiment}
\label{sec:supp-cross-verifier}

This experiment uses
\texttt{Qwen3-4B-Instruct-2507} and
\texttt{Meta-Llama-3.1-8B-Instruct}.  Each backbone is LoRA-tuned under two
complementary mixed-arm assignments and three random seeds, yielding
\(2\times2\times3=12\) runs.  In one assignment, each training claim is
assigned to exactly one evidence arm; the complementary assignment reverses
the arm for every claim.  Each assignment contains 5,025 EviTrace and 5,025
retrieval-ordered examples, balanced within
gold-label \(\times\) single-/multi-atom strata.  Checkpoints are selected
by the mean validation Macro-F1 over the two arms.

Both backbones use their native chat templates and the same method-blind
LIAR-6 label prompt.  The frozen recipe is LoRA rank 16, alpha 32, dropout
0.05, attention and MLP projection targets, BF16 ZeRO-2, effective batch
size 16, learning rate \(2\times10^{-5}\), 12 epochs, and maximum length
2,048 tokens.

The \emph{Main} comparison uses 1,250 test claims.  It contrasts the final
verifier-visible EviTrace sequence with a source-score retrieval-only
top-\(K_i\) sequence drawn from the same Atom--Union pool.  The evidence
count is equal within each claim, but token counts are not forced to be
identical.  The \emph{Order-only} comparison fixes the selected evidence
identities, texts, and count and changes only their order: EviTrace order
versus descending source-score order.  The 98 claims for which the two
orders are identical are excluded by this comparison's fixed eligibility
rule, leaving 1,152 claims.

Predictions are first averaged over assignment and seed within a backbone
and are then combined with equal backbone weight.  Claims, not the 12
repeated predictions, are the independent units.  A
gold-label-stratified, claim-clustered bootstrap with 10,000 replicates
provides the 95\% confidence interval for the paired Macro-F1 difference.
The randomization test swaps the two arm predictions independently at the
claim level 100,000 times.  The Main and Order-only hypotheses form one
family and their two \(p\)-values are Holm-adjusted.

For descriptive paired accuracy, a claim is a win when EviTrace alone is
correct, a loss when the control alone is correct, and a tie when both are
correct or both are incorrect.  The conditional win rate is
\begin{equation}
\operatorname{CWR} =
\frac{N(\text{EviTrace correct, control wrong})}
{N(\text{exactly one arm correct})}.
\end{equation}
Main obtains a \(+1.29\)-point Macro-F1 difference
(95\% CI \([0.26,2.35]\), Holm-adjusted \(p=.029\)) and a 53.95\%
conditional win rate.  Order-only obtains \(+0.67\) points
(95\% CI \([-0.04,1.39]\), Holm-adjusted \(p=.065\)) and a 52.97\%
conditional win rate.  The latter interval crosses zero and is interpreted
as suggestive rather than conclusive evidence of an ordering-only effect.

\subsection{Evidence-Capacity Sensitivity for RQ3}
\label{sec:supp-rq3}

All capacity conditions reuse the same ordered trace and differ only in the
prefix made visible to the verifier.  A fixed-\(k\) policy retains the first
\(k\) complete evidence units.  A minmax\((k_{\min},k_{\max})\) policy
shows at least \(k_{\min}\) units, then stops at the first stored
\texttt{target\_resolved} flag---derived from the replayed resolved-atom
rate with \(\rho_{\mathrm{target}}=1.0\)---and otherwise continues up to
\(k_{\max}\).  A budget-\(B\) policy appends complete units while its
capacity guard remains satisfied.  Pool exhaustion may produce an
underfilled prefix and is recorded separately from a resolution stop.  The
later tokenizer guard normally removes complete tail units; only its
single-unit emergency fallback may truncate within evidence text, as
specified in Section~\ref{sec:supp-truncation}.

\begin{table}[t]
\centering
\scriptsize
\setlength{\tabcolsep}{3pt}
\begin{tabular}{llrr}
\toprule
\textbf{Family} & \textbf{Setting} & \textbf{Macro-F1}
& \textbf{Mean toks.} \\
\midrule
fixed & \(k=3\) & 34.12 & 436.33 \\
fixed & \(k=5\) & 33.91 & 603.45 \\
fixed & \(k=7\) & 34.03 & 751.29 \\
fixed & \(k=9\) & 34.81 & 867.65 \\
\midrule
minmax & \((3,8)\)  & 33.84 & 540.19 \\
minmax & \((3,10)\) & 33.86 & 562.51 \\
minmax & \((5,10)\) & 36.66 & 677.90 \\
minmax & \((5,12)\) & 35.22 & 686.65 \\
minmax & \((7,12)\) & 34.50 & 798.46 \\
\midrule
budget & \(B=512\)  & 32.88 & 461.36 \\
budget & \(B=768\)  & 34.45 & 710.40 \\
budget & \(B=1024\) & 35.03 & 949.29 \\
\bottomrule
\end{tabular}
\caption{Complete LIAR-RAW capacity grid underlying the RQ3 figure.
Prompt-token means are computed after final rendering with the verifier
tokenizer.}
\label{tab:supp-capacity-grid}
\end{table}

The policy family, parameter values, trace fingerprint, tokenizer revision,
and prompt template are stored with every point.  The capacity plot reports
all evaluated settings rather than only the best setting.  Its connecting
lines are visual guides, not fitted learning curves, and the table contains
point estimates rather than a claim of monotonicity.

\placeholder{State the validation criterion used to select
minmax(5,10) as the primary setting and confirm that no test metric was used
for policy selection.}

\placeholder{Add the median, 95th percentile, maximum, actual
visible-\(K\) distribution, and truncation/underfill rates for every
capacity condition.}

\subsection{Human Audit Protocol and Metrics}
\label{sec:supp-human}

\subsubsection{Sampling and annotation procedure}

The atomization audit draws a frozen validation-set sample from LIAR-RAW and
RAWFC.  Sampling combines a uniform component with a
difficulty-prioritized component containing long claims or claims with
negation, comparison, numerical, date, or extremum cues.  Exact item
identifiers and the sampling script are retained in the audit manifest;
sample totals and summary outcomes remain in the main paper.

Two annotators independently assess atom faithfulness, atomicity, and
claim-level completeness.  Faithfulness asks whether an atom is entailed by
the original claim without adding information.  Atomicity asks whether it
is a minimal independently verifiable proposition.  Completeness records
how many independently verifiable propositions in the claim are omitted by
the atom set.  Inter-annotator agreement is computed on the two initial
annotations before adjudication; adjudication is used only to construct the
final audit label.

The Evidence Map audit draws candidate--atom pairs from both datasets and
follows the natural relation distribution rather than balancing classes.
Two annotators independently assign relation, directness, and self-reported
confidence; both annotations are retained as empirical human references.
The annotation task is a local text-pair judgment and does not ask
annotators to search the Web or determine the claim's overall verdict.

\placeholder{Describe annotator educational/professional
background and relevant NLP or fact-checking experience.}

\placeholder{State whether each annotator and adjudicator was
an author, and describe any prior exposure to the sampled items.}

\placeholder{Report the training/calibration procedure,
guideline version, calibration set size, acceptance criterion, and whether
running agreement checks were performed.}

\placeholder{State the implemented blinding: mutual annotator
blinding; visibility of LLM labels, model confidence, and veracity labels;
and the information shown to the adjudicator.}

\placeholder{Report compensation, payment basis, approximate
workload and completion time per annotator, and recruitment procedure.}

\placeholder{Provide the institutional ethics/IRB
approval, exemption, or non-human-subjects determination and the consent
procedure.}

\subsubsection{Audit label spaces}

Atom faithfulness and atomicity are binary.  Completeness uses
\(\{0,1,2,3+\}\) missed propositions; the reported complete-coverage
statistic is the binary event \(\textit{missed}=0\).
The implementation-level human relation inventory is
\[
\begin{aligned}
\{&\textit{support},\textit{refute},\textit{qualify},\textit{mixed},\\
&\textit{insufficient},\textit{background},\textit{irrelevant}\},
\end{aligned}
\]
and directness is
\(\{\textit{direct},\textit{partial},\textit{context},\textit{none}\}\).
The compact relation display in the main paper is an abstract grouping for
explaining state transitions, not the stored annotation schema.  Human
agreement and LLM--human alignment are computed in the implementation-level
label spaces above.

\subsubsection{Atomization metrics}

For \(N\) double-annotated units, exact agreement is
\[
A_{\mathrm{exact}}=\frac{1}{N}\sum_{i=1}^{N}
\mathbb{1}[a_i=b_i].
\]
Cohen's kappa is
\[
\kappa=\frac{p_o-p_e}{1-p_e},
\]
where \(p_o\) is observed agreement and \(p_e\) is agreement expected from
the two annotators' marginal label frequencies.  Because the atom-quality
labels are highly imbalanced, the paper also reports Gwet's AC1.  For \(q\)
categories with mean marginal proportion \(\pi_k\),
\[
\begin{aligned}
P_e^{\mathrm{AC1}}
&=\frac{1}{q-1}\sum_{k=1}^{q}\pi_k(1-\pi_k),\\
\mathrm{AC1}
&=\frac{p_o-P_e^{\mathrm{AC1}}}
{1-P_e^{\mathrm{AC1}}}.
\end{aligned}
\]

After adjudication, the pass rate is the fraction of units whose final
label passes the named criterion.  A claim receives the strict all-pass
label only when completeness has zero omissions and every atom is both
faithful and atomic.  Confidence intervals for final pass rates use 5,000
dataset-stratified, claim-clustered bootstrap replicates; all atoms from one
claim stay in the same resample.  Pre-adjudication agreement and
post-adjudication pass rate are reported separately because they answer
different questions.

\subsubsection{Evidence Map agreement and alignment}

Relation reliability uses exact agreement and Cohen's \(\kappa\).
Directness assigns ordinal scores
\(\textit{none}=0\), \(\textit{context}=1\),
\(\textit{partial}=2\), and \(\textit{direct}=3\), and reports Spearman's
rank correlation.  Within-one agreement is
\[
A_{\leq1}=\frac{1}{N}\sum_i
\mathbb{1}\!\left[|d_i^{(A)}-d_i^{(B)}|\leq1\right].
\]
For an LLM label \(\hat y_i\) and human labels
\(y_i^{(A)},y_i^{(B)}\), average exact alignment is
\[
A_{\mathrm{avg}}=
\frac{1}{2N}\sum_i
\left(\mathbb{1}[\hat y_i=y_i^{(A)}]
+      \mathbb{1}[\hat y_i=y_i^{(B)}]\right),
\]
whereas either-annotator alignment is
\[
A_{\mathrm{either}}=
\frac{1}{N}\sum_i
\mathbb{1}[\hat y_i\in\{y_i^{(A)},y_i^{(B)}\}].
\]
The corresponding directness measures use the same definitions, with the
within-one tolerance applied to the ordinal scores.

For calibration, let \(c_i\in[0,1]\) be the LLM confidence and define the
soft human-match target
\[
z_i=\tfrac12\left(
\mathbb{1}[\hat y_i=y_i^{(A)}]
+\mathbb{1}[\hat y_i=y_i^{(B)}]\right).
\]
Using five equal-width confidence bins \(\mathcal B_b\), define
\[
\begin{aligned}
\bar c_b&=\frac{1}{|\mathcal B_b|}
\sum_{i\in\mathcal B_b}c_i,\\
\bar z_b&=\frac{1}{|\mathcal B_b|}
\sum_{i\in\mathcal B_b}z_i.
\end{aligned}
\]
Expected calibration error~\citep{Guo2017Calibration} and Brier score are
\begin{align}
\mathrm{ECE}
&=\sum_{b=1}^{5}\frac{|\mathcal B_b|}{N}
\left|\bar c_b-\bar z_b\right|,\\
\mathrm{Brier}
&=\frac{1}{N}\sum_{i=1}^{N}(c_i-z_i)^2.
\end{align}
Human confidence values are annotator self-reports, not factual gold
probabilities.  The summary values remain in the main paper; this supplement
fixes the metric definitions used to obtain them.

\subsection{Case-Study Reading Guide}
\label{sec:supp-case}

The main paper's case-study figure visualizes one complete, source-linked
trace.
The top block contains the original claim and its atom identifiers.  Each
trace row then records the selected evidence identifier and text, its
source report, the mapped relation/directness/confidence, its retrieval
score and state-conditioned utility, and the atom-state transition induced
after appending the unit.  The bottom block contains the stop reason, the
exact verifier-visible ordered prefix, the allowed label tokens, the
predicted label, and the dataset label.  Candidate identities remain stable
across the retrieval pool, map, trace, and rendered prompt, so the displayed
sequence can be replayed from the released artifacts.

Candidates that are retrieval-relevant but add no prefix-dependent utility
remain in the candidate ledger with a ``not selected'' status.  This makes
the example useful for inspecting what the selector exposed to the
verifier and how its explicit atom states changed.  It does not reveal the
verifier's latent chain of thought, and the displayed state trace is not
claimed to be a causal explanation of the verifier prediction.

\placeholder{State the case-selection rule (random,
pre-registered diagnostic, or illustrative), dataset/split/event ID, gold
and predicted labels, and whether the example was selected before viewing
the final prediction.}

\placeholder{Add the stable report/chunk identifiers and
source-link availability statement for every evidence unit in the figure.}

\placeholder{Add at least one incorrect-prediction or
conflicting-evidence example, or state explicitly that Figure 4 is a single
illustrative example rather than a representative qualitative sample.}

\subsection{Limitations}
\label{sec:supp-limitations}

\paragraph{Scope and external validity.}
The experiments cover English political claims, Snopes claims, and
scientific claims, with the controlled RQ2/RQ3 analyses concentrated on
LIAR-RAW and one primary 8B verifier.  The cross-verifier experiment adds
a 4B and an 8B verifier but does not establish invariance across model scales,
languages, retrieval corpora, or time periods.

\paragraph{Upstream dependence.}
EviTrace cannot recover evidence absent from its candidate pool.
Atomization errors, omitted qualifiers, report parsing errors, and
LLM-generated Evidence Map errors can propagate through state updates and
stopping.  Human auditing supports the quality of the sampled structures
but does not establish that every generated annotation is correct.  A
claim-linked report is also not necessarily a credible source: reports may
repeat a false statement, omit context, or duplicate one another.

\paragraph{Selector and capacity assumptions.}
The autoregressive selector is greedy and is not guaranteed to find a
globally optimal sequence.  Its stopping rule treats map-derived atom
resolution as an operational proxy for evidence sufficiency, which need
not coincide with logical sufficiency or a human fact-checker's judgment.
Changes to map fields may affect both selected content and visible
capacity.  Conflicting evidence is represented in the explicit atom state,
but the final resolution of that conflict is delegated to the verifier.

\paragraph{Interpretation boundary.}
The trace is an inspectable record of retrieved units, selection order, and
explicit state transitions.  It is not a faithful transcript of the
verifier's hidden reasoning and does not, by itself, prove that a selected
unit causally determined the prediction.  The single case study illustrates
artifact-level auditability, not population-level interpretability.

\paragraph{Statistical and operational uncertainty.}
\placeholder{State which reported cells are single runs and
which aggregate multiple seeds, then describe the resulting uncertainty
and any unmeasured API nondeterminism.}
The RAWFC test set contains only 200 claims, so small point differences
should be interpreted cautiously.  LLM annotation also introduces
provider dependence, cost, possible model updates behind an API name, and
failure modes that may not be reproduced without the frozen response
cache.

\paragraph{Benchmark exposure and temporal drift.}
The external annotation models and pretrained verifier backbones may have
encountered benchmark claims or related fact checks during pretraining; this
cannot be ruled out from public model documentation alone.  Moreover,
evidence availability and the truth status of time-sensitive claims can
change after dataset construction.  Results should therefore be interpreted
as benchmark measurements at the frozen data and model snapshots, not as
guarantees for future claims.

\subsection{Ethical Considerations}
\label{sec:supp-ethics}

\paragraph{Human annotation.}
Annotators evaluated text pairs and generated structures; they did not
make deployment decisions about people or organizations.
\placeholder{Insert the ethics/IRB determination, informed
consent language, recruitment, compensation, workload, annotator
background, author/non-author status, and data-retention policy.}
Because political and fact-checking reports can contain offensive language
or allegations about identifiable people, the final version should also
state what content warning, withdrawal option, and support procedure were
provided.

\paragraph{Potential misuse and human oversight.}
Incorrect fact-checking labels can cause reputational, political, or
informational harm.  EviTrace should not be used as an autonomous arbiter,
content-removal system, or substitute for professional review in
high-stakes settings.  A deployment should retain source links, expose
uncertainty and conflicting evidence, permit correction and appeal, and
require a human reviewer before consequential action.

\paragraph{Bias and representativeness.}
The datasets reflect the editorial practices, topic distributions, and
source coverage of their creators.  Political ideology, demographic
groups, less-covered regions, and non-English claims may therefore receive
unequal evidence coverage or error rates.  A source's inclusion in a
claim-linked bundle must not be interpreted as a credibility endorsement.
Subgroup and temporal audits are required before use outside the benchmark
setting.

\paragraph{Privacy, copyright, and licensing.}
Raw reports can contain names, quotations, and copyrighted text.  The
artifact should distribute identifiers, manifests, and reconstruction
instructions when redistribution rights do not permit republishing the
full report text.  Dataset, model, embedding-model, and API licenses and
terms of service must be listed, and any personal data should be handled
under the source dataset's research-use conditions.
\placeholder{Confirm the redistribution decision and license
for each dataset, model adapter, prompt/response cache, and annotation
export.}

\paragraph{External-service data handling.}
Atomization and Evidence Map construction transmit the claim and the
minimum required evidence text to an external model endpoint.  A production
deployment should avoid sending secrets or unnecessary personal data and
should provide a local-processing alternative where confidentiality
requires it.
\placeholder{Identify the API provider and region; state the
applicable terms, logging/retention and model-training policy, deletion
mechanism, access controls, encryption, and any institutional data-processing
approval.}

\paragraph{Compute and environmental cost.}
Training and retrieval used NVIDIA L20 accelerators, and preprocessing uses
an external LLM API.
\placeholder{Clarify node/GPU count, L20 memory, CPU/RAM,
software versions, wall-clock time, GPU-hours, API request/token counts,
monetary cost, and any energy/carbon estimate or mitigation.}

\subsection{Reproducibility Artifact Checklist}
\label{sec:supp-artifacts}

Tables~\ref{tab:supp-artifact-checklist-core}
and~\ref{tab:supp-artifact-checklist} list the minimum reconstruction
package.  The split presentation is only for two-column layout; together
the two tables form one checklist.

\begin{table*}[t]
\centering
\scriptsize
\setlength{\tabcolsep}{3.1pt}
\begin{tabularx}{\textwidth}{@{}p{0.13\textwidth}X X
p{0.12\textwidth}@{}}
\toprule
\textbf{Artifact} & \textbf{Required contents}
& \textbf{Integrity/provenance field} & \textbf{Release status} \\
\midrule
Data manifest &
Source URL, release date, split counts, class counts, licenses, raw-file
paths &
SHA-256 of every downloaded file &
\placeholder{TODO} \\
Preprocessing code &
Text cleaning, sentence splitting, chunking, report materialization,
label normalization &
Git commit and resolved configuration &
\placeholder{TODO} \\
Prompt/schema bundle &
Atomization and map prompts, JSON schemas, parsers, retry and fallback
rules &
Prompt/schema hashes and API model/date &
\placeholder{TODO} \\
LLM annotation cache &
Frozen requests, responses, parsed records, errors, and retry metadata &
Per-file hash and provider request metadata &
\placeholder{TODO} \\
Retrieval bundle &
Embedder revision, chunk manifest, route results, deduplication, pool
fingerprints &
Model revision and cache fingerprints &
\placeholder{TODO} \\
Scorer bundle &
Preference pairs, feature definitions, resolved training config, weights,
and supervision audit &
Weight hash, seed, and input-manifest hashes &
\placeholder{TODO} \\
Trace bundle &
Candidate ledger, selected order, per-step score/state, stop reason, and
visible-prefix fingerprint &
Stable event/evidence IDs and JSONL hashes &
\placeholder{TODO} \\
\bottomrule
\end{tabularx}
\caption{Core data, preprocessing, mapping, scorer, and trace artifacts
required for reconstruction.}
\label{tab:supp-artifact-checklist-core}
\end{table*}

\begin{table*}[t]
\centering
\scriptsize
\setlength{\tabcolsep}{3.1pt}
\begin{tabularx}{\textwidth}{@{}p{0.13\textwidth}X X
p{0.12\textwidth}@{}}
\toprule
\textbf{Artifact} & \textbf{Required contents}
& \textbf{Integrity/provenance field} & \textbf{Release status} \\
\midrule
Verifier bundle &
Base-model revision, LoRA adapter, tokenizer/chat template, resolved
optimizer config, checkpoint choice &
Model/adapter/tokenizer hashes and seed &
\placeholder{TODO} \\
Evaluation bundle &
Claim-level predictions, label maps, metric implementation, confusion
matrices, and raw metric JSON &
Prediction and script hashes &
\placeholder{TODO} \\
RQ2 statistics &
The 12 cross-verifier run outputs, assignment manifests, bootstrap and
arm-swap scripts &
Seeds, claim universe hash, Holm family definition &
\placeholder{TODO} \\
RQ3 grid &
All 12 rendered builds, prompt-token counts, visible-\(K\), truncation and
underfill diagnostics &
Shared-trace and tokenizer fingerprints &
\placeholder{TODO} \\
Human audit &
Sampling script/stats, task files without revealed LLM labels, guideline,
raw double annotations, adjudication, metric script &
Task/guideline/result hashes and anonymized annotator IDs &
\placeholder{TODO} \\
Case studies &
Case-selection rule, complete candidate ledger, source identifiers,
rendered trace, prediction &
Event ID and source/trace fingerprints &
\placeholder{TODO} \\
Environment and cost &
Hardware topology, OS/driver/CUDA, library lockfile, wall time, GPU-hours,
API usage &
Machine-readable environment manifest &
\placeholder{TODO} \\
\bottomrule
\end{tabularx}
\caption{Verifier, evaluation, audit, case-study, and environment artifacts
required for reconstruction.}
\label{tab:supp-artifact-checklist}
\end{table*}

Every released manifest should record the creating command, UTC timestamp,
code commit, resolved configuration, upstream input hashes, output hashes,
and completion status.  A mutable \texttt{latest} path is not an
authoritative citation; tables and figures should point to immutable run
identifiers.  The artifact landing page, archival DOI, and release license
are:
\placeholder{insert repository URL, archival DOI, release
date, and license.}
